% RESEARCH_DRAFT_SHA256: b9a620e0d72e85182b13ffca41be61dfa332d9106b772932167979c3ee085803
% Layout and source-reference wrappers expanded for the web renderer.
\documentclass[11pt]{article}
\usepackage[margin=25mm]{geometry}
\usepackage{amsthm}
\usepackage{fontspec,unicode-math}
\setmainfont{Latin Modern Roman}
\setmathfont{Latin Modern Math}
\usepackage{amsmath,amssymb,mathtools}
\usepackage{xurl,hyperref}
\hypersetup{colorlinks=true,urlcolor=blue,linkcolor=blue}
\setcounter{secnumdepth}{0}
\setlength{\parindent}{0pt}
\setlength{\parskip}{5pt}
\setlength{\emergencystretch}{3em}
\newcommand{\sref}[2]{\hyperlink{source:#1:#2}{[S#2]}}
\newcommand{\cataloguescope}[1]{\paragraph{Source question.}#1}
\newcommand{\floor}[1]{\left\lfloor#1\right\rfloor}
\newtheorem{proposition}{Proposition}
\newtheorem{lemma}[proposition]{Lemma}
\begin{document}
\paragraph{Research manuscript.} OpenAI \texttt{gpt-6-astra}, reasoning \texttt{ultra}; 11 October 2026.
\paragraph{Preserved source body.} The complete catalogue mathematical text and sources follow with layout and source links expanded, including source qualifications and its historical preparation notes. New research begins at the section headed Research attempt.

\section{Sharp Smoothed Complexity of Simplex Pivot Rules}\label{statement}

\subsection{Definitions and mathematical statement}

\paragraph{Shared notation.}Standard number systems are denoted by $\mathbb N,\mathbb Z,\mathbb Q,\mathbb R,\mathbb C$; index ranges, computational models and specialized notation follow the cited problem source and the definitions below.


Smoothed analysis studies expected performance after random input perturbation. A projected polyhedron is the image under a linear map; its edges encode steps in the relevant shadow-vertex simplex analysis.

A linear program optimizes $c^Tx$ over a polyhedron $\{x:Ax\leq b\}$. In the simplex method a pivot moves between adjacent basic feasible solutions; a pivot rule chooses the next such move.

A polyhedron is an intersection of finitely many closed affine half-spaces; a polytope is a bounded polyhedron. Its faces are its intersections with supporting hyperplanes.
\paragraph{Statement recorded in the supplied source.}
For smoothed analysis of the simplex method, can the expected number of edges of a projected polyhedron be bounded sharply by $\Theta(d^{3/2}\log n)$ as $\sigma\to\infty$, with a lower bound depending on $\sigma$? What bounds hold if only a random $\varepsilon$-fraction of the entries of $(A,b)$ is perturbed, or if other pivot rules are used? \href{https://doi.org/10.4171/owr/2018/54}{[S2]}
\paragraph{Scope qualification.}
The normalization of $A,b$, Gaussian variance $\sigma$, projection, and permitted pivot rules must be taken from the report.
\paragraph{Source status.}
The supplied catalogue labels this record open. Its source-status metadata was checked against the supplied snapshot on 2026-10-11. This import does not independently establish current open status or mathematical correctness. The catalogue literature assessment is dated 2026-08-21.
\subsection{Short English statement}
Sharpen the stated simplex smoothed bounds and extend them to sparse perturbations and other pivot rules.
\subsection{Sources}
\begin{itemize}
\item[S1] Ulam AI and the UnsolvedMath Contributors, UnsolvedMath: A Curated Collection of Open Mathematics Problems, catalogue entry. Catalogue curation and metadata: CC BY 4.0. The supplied source question is retained; definitions and the English summary are editorial additions.
\url{https://huggingface.co/datasets/ulamai/UnsolvedMath}.
\item[S2] Convex Geometry and its Applications (2019), Oberwolfach Reports, DOI 10.4171/owr/2018/54, pp. 3217--3281
\url{https://doi.org/10.4171/owr/2018/54}.
\item[S3] Dadush--Huiberts
\url{https://doi.org/10.1137/18M1197205}.
\item[S4] Huiberts--Lee--Zhang
\url{https://arxiv.org/abs/2211.11860}.
\textit{Underlying Oberwolfach Reports material retains its source-specific licence, including CC BY-SA 4.0 where applicable; see THIRD\_PARTY\_NOTICES.md.}
\end{itemize}
\paragraph{Preparation and review.}
Original question attributed to the cited source and UnsolvedMath contributors. Codex (GPT-6; exact serving revision not exposed) assisted with definitions, English summaries, source inspection and formatting. No new human mathematical authorship or independent proof verification is asserted. Catalogue identity checked against the complete supplied registry and existing repository statements; known flagged repetitions and obvious identical conjectures excluded. Full mathematical-equivalence/literature review has not been performed. Statement preparation and final review followed contribution prompts C and D; automated checks verify structure and identity, not open status.
\label{end:statement}

\section{Research attempt}
\subsection{Literature, source context and scope}
The full original catalogue body and source qualifications are preserved above. The relevant primary workshop passages were read in the cached publisher report. Historical claims and subsequent attributed results are distinguished from the mathematical work below. A restricted example, standard lemma or known theorem is not counted as new progress on the entire original target. Missing assumptions and repeated or rhetorical extracts are not silently turned into different solved problems.

\subsection{A source discrepancy and a sparse-perturbation obstruction}
The preserved catalogue displays $\Theta(d^{3/2}\log n)$. Inspection of the actual primary PDF image on printed page 3230 shows $\Theta(d^{3/2}\sqrt{\ln n})$ in its fixed-$d$, centered Gaussian comparison. These are different targets; the original body is retained unchanged and neither sharp bound is proved here. The same primary page asks about perturbing only a random fraction of entries. We test whether full-Gaussian almost-sure nonsingularity can simply be carried over to independent sparse perturbations.
Let $M=\begin{pmatrix}1&1\\1&1\end{pmatrix}$. Independently for each of its four entries, with probability $\varepsilon\in[0,1]$ add a centered Gaussian of variance $\sigma^2>0$, and otherwise leave it unchanged. This Bernoulli selection model is one explicit meaning of sparse perturbation; it is not sampling an exactly fixed-size subset.
\begin{proposition}
The perturbed matrix is singular with probability exactly $(1-\varepsilon)^4$, for every $\sigma>0$. In particular this singular mass does not disappear as $\sigma$ tends to infinity when $\varepsilon<1$.
\end{proposition}
\begin{proof}
If no entry is selected, the original singular matrix remains, an event of probability $(1-\varepsilon)^4$. Conditional on any nonempty selection, the determinant is a nonzero polynomial in the selected independent Gaussian increments: set all increments but one selected entry to zero, and that entry changes the determinant linearly with coefficient 1 or $-1$. A nonzero polynomial evaluated at variables with a joint density vanishes with probability zero. To see this last fact, induct on the number of variables. In one variable there are finitely many roots. In more variables, write the polynomial as a polynomial in the last variable with coefficient polynomials in the others. At least one coefficient is nonzero; by induction the event that it vanishes has measure zero. Outside that event the last-variable polynomial has finitely many roots, and integration of the conditional zero probabilities finishes the argument. Every nonempty selection therefore contributes zero singular probability, proving the formula.
\end{proof}
\subsection{Verification and remaining gaps}
All 16 selection masks were enumerated. For each nonempty mask, changing one selected entry by 1 while leaving the others at zero gave determinant $\pm1$, verifying nonidentity of its determinant polynomial. The atom probability was evaluated exactly at rational selection probabilities. This blocks an argument that assumes the sparse-perturbation distribution is fully absolutely continuous or automatically removes every degeneracy. It does not imply infinite simplex runtime: degeneracy can be handled by appropriate pivot conventions, and expected projected-edge bounds need a separate analysis. No sharp shadow-size estimate, sigma-dependent lower bound, or guarantee for other pivot rules is obtained.

\section{References and source attribution}
Primary workshop report: \url{https://doi.org/10.4171/owr/2018/54}. Relevant passages read: 2018/54 pp. 3229--3230, Gaussian normalization, projected-edge theorem and all three open-question clauses.

\url{https://ems.press/content/serial-article-files/46776}. Primary PDF p.3230 image directly inspected using local PDFKit rendering after web screenshot cache failure; confirmed sqrt(ln n) in the comparison bound.

\url{https://arxiv.org/abs/2211.11860}. Opened primary later-paper abstract page on upper/lower smoothed simplex bounds; no claim of having independently verified its full parameter theorems.

All retrievals were checked on 11 October 2026. Original question attribution remains in the preserved Sources section. External theorems are attributed and not claimed as new discoveries or fully verified merely by citation.


\paragraph{Integration review.} Codex (GPT-6; exact serving revision not exposed) reviewed the written arguments and expanded layout and source-reference wrappers. Repository metadata and the production web renderer were checked using Python 3.13 and JavaScript with MathJax 3.2.2. This model-assisted review is not independent human proof verification; no PDF compilation or formal proof checking is asserted.

\section{Completion Estimate}
\noindent\textbf{COMPLETION: 0\%}
No substantiated new progress toward the entire original target is claimed. The worked results and failed route are delimited above; a known result, illustrative model, or auxiliary obstruction is not counted as original-target completion.

\section{Final status and attribution}
\textbf{Attempt status: unresolved.} This describes this research attempt, not a new ruling on the source catalogue's status or on external literature claims.
The mathematics and this self-audit were generated by OpenAI \texttt{gpt-6-astra}, reasoning \texttt{ultra}, on 11 October 2026. Codex (GPT-6; exact serving revision not exposed) orchestrated the batch, reviewed the written mathematical arguments, and contributed formatting/repository review. This additional model review is not independent human verification. No human mathematical authorship, independent human proof checking, or formal verification is claimed. The historical preparation paragraphs in the reproduced catalogue body concern that body only.

Contribution prompt A was applied to the complete definition and supplied source record, followed by the adversarial checks of prompt B. Local reproduction files are \texttt{batch022/checks.py}, \texttt{computations.json}, and the per-ID audit JSON. They were not public links at preparation time. Numerical checks are finite and do not prove unrestricted results. No LaTeX engine was found locally; compilation was not performed by this research agent.

\end{document}
